\begingroup
\scriptsize
\setlength{\LTcapwidth}{0.9\textwidth}
\begin{longtable}[width=0.9\textwidth]{p{0.2\textwidth}p{0.3\textwidth}p{0.4\textwidth}}
\caption{Taxonomy of pragmatic aspects of narrative intent and reception of storytelling on social media}
\label{tab:taxonomy} \\
\toprule
\rowcolor[gray]{0.75}
\multicolumn{1}{p{0.2\textwidth}}{\textbf{Sub-dimension}} & \multicolumn{1}{p{0.3\textwidth}}{\textbf{Definition}} & \multicolumn{1}{p{0.4\textwidth}}{\textbf{Example}} \\
\midrule
\endfirsthead
\multicolumn{3}{p{0.9\textwidth}}{\tablename\ \thetable{} -- continued from previous page} \\
\toprule
\rowcolor[gray]{0.75}
\multicolumn{1}{p{0.2\textwidth}}{\textbf{Sub-dimension}} & \multicolumn{1}{p{0.3\textwidth}}{\textbf{Definition}} & \multicolumn{1}{p{0.4\textwidth}}{\textbf{Example}} \\
\midrule
\endhead
\midrule
\multicolumn{3}{r}{Continued on next page} \\
\endfoot
\bottomrule
\endlastfoot
\rowcolor{SkyBlue!15}
\multicolumn{3}{p{0.953\textwidth}}{
\textbf{Dimension: Overall Goal}\newline
\textbf{Template}: \textit{Many readers from this subreddit would think that the author's overall goal in posting/commenting this text was to \{\{SHORT VERB PHRASE DESCRIBING OVERALL GOAL\}\}.}\newline
\textbf{Summary}: Some common overall goals for posting on social media include:\par - request or provide facts or info about approaches, strategies, etc\par - seek or provide emotional support\par - express/reinforce one’s identity, values, or accomplishments\par - solicit or share experiences to gain or share insights\par - argue/advocate for a viewpoint to persuade others\par - share an enjoyable or funny post
} \\
\midrule
request\_info\_support & request factual info or advice about approaches, strategies, etc. & The author recounts a series of troubleshooting steps they've already followed to solve a technical issue, then poses a technical question. \\
provide\_info\_support & provide facts or advice about approaches, strategies, etc. & The author reports a newsworthy event relevant to the community. \\
request\_emotional\_support & express emotions or characterize a situation to elicit others' comfort, understanding, or empathy. & The author describes a series of unlucky events with their family and reports feeling misunderstood. \\
provide\_emotional\_support & provide emotional support to someone by acknowledging their identity, values, or accomplishments; or offering emotional comfort, understanding, or empathy. & The author justifies how someone who (presumably) posted a text earlier in the conversation reacted in some situation. \\
affirm\_identity\_self & reinforce or assert their own identity, values, or accomplishments. & The author asserts that they are a good student. \\
request\_experiential\_accounts & solicit stories or experience from others to gain perspective or insight. & The author describes a strange experience and then asks if anyone else has experienced something similar. \\
provide\_experiential\_accounts & share a personal story or experience to inform or engage others. & The author shares a story about their immigration experience to show that it was different than how someone who (presumably) posted earlier characterized immigration. \\
persuade\_debate & advocate for a viewpoint or make an argument to convince others. & The author asserts that a sports team made a bad trade based on the poor track record of a player. \\
entertain & share an enjoyable or funny post. & The author repeats a famous joke from a comedian. \\
\midrule
\rowcolor{SkyBlue!15}
\multicolumn{3}{p{0.953\textwidth}}{
\textbf{Dimension: Narrative Intent}\newline
\textbf{Template}: \textit{Many readers from this subreddit would think that the author told the story in their post/comment to \{\{SHORT VERB PHRASE DESCRIBING NARRATIVE INTENT\}\}.}\newline
\textbf{Summary}: Some common reasons people tell a story in a social media post include:\par - express/demonstrate one’s identity or beliefs\par - teach a life lesson\par - justify or challenge a belief or social norm\par - entertain\par - release pent up emotions (vent)\par - convey need for emotional support\par - convey a similar experience\par - correct misunderstandings or add missing details associated with an event/situation
} \\
\midrule
show\_identity & demonstrate an aspect of their identity in action, by example. & The author shows that they can't be bought by telling a story about how they turned down a lucrative job offer that did not align with their values. \\
justify\_challenge\_offer\_belief\_norm & explain how they came to hold or question a belief or social norm, or to reinforce/demonstrate why it is correct, beneficial, misguided, or harmful; teach a life lesson or influence the behaviors or attitudes of readers & The author told a story about how they came to appreciate the politeness of small talk in the South in the U.S. \\
entertain & share an enjoyable or funny story. & The author tells a story about a dream where they were giving a presentation in their underwear. \\
release\_pent\_up\_emotions & express themself as a means of emotional release or sensemaking. & The author recounts being totally caught off guard in a Zoom call where they were fired unceremoniously. \\
convey\_emotional\_support\_need & describe events that have left them in an unfortunate or unresolved state, in need of emotional support. & The author tells a story about how they have become lonely over the last few months because there friends have stopped reaching out. \\
convey\_similar\_experience & share a story similar to another story or related to a situation under discussion. & The author tells a story about their experience with graduate school admissions in the context of a discussion about different people's similar experiences. \\
clarify\_what\_transpired & correct misunderstandings or add missing details associated with an event or situation under discussion. & The author points out that a discussion has overlooked several important factors or events that are necessary to understand a situation under discussion. \\
\midrule
\rowcolor{SkyBlue!15}
\multicolumn{3}{p{0.953\textwidth}}{
\textbf{Dimension: Author Emotional Response}\newline
\textbf{Template}: \textit{Many readers from this subreddit would think that telling the story in their post/comment would cause the author of the post/comment to feel \{\{SHORT NOUN PHRASE DESCRIBING EMOTION\}\}.}\newline
\textbf{Summary}: Some feelings an author could experience after telling a story in a social media post include: fear, guilt, anger, sadness, disgust, envy, joy, pride, relief, hope, compassion, appreciation, and connection.
} \\
\midrule
fear & a response to perceived danger or threat. & The author tells a story about how they heard banging on their door in the middle of the night. \\
guilt & remorse for violating personal or social standards. & The author tells a story about how they have been keeping the fact that they dropped out of college a secret from their parents. \\
anger & a strong reaction to perceived harm, injustice, or frustration. & The author tells a story about how their dog was attacked at a park by another dog whose owner did not use a leash. \\
sadness & a sense of loss, disappointment, or helplessness. & The author tells a story about how they were rejected from their dream school. \\
disgust & a reaction of revulsion to something perceived as offensive or repellent. & The author tells a story about how they saw someone urinating on the subway. \\
envy & a desire for something others have, coupled with resentment. & The author tells a story about how their coworker was promoted over them despite not being as good an employee as the author.  \\
joy & a state of happiness and contentment. & The author tells a story about their wedding night that they remember fondly. \\
pride & a sense of satisfaction from achievements or qualities. & The author tells a story about their coming out process and the ways their life has changed for the better since coming out. \\
relief & a release from stress or tension after resolving a concern. & The author tells a story about narrowly avoiding being caught skipping the last hour of their work shift. \\
hope & an optimistic expectation for a positive outcome. & The author tells a story about their persistence trying to learn the piano and their excitement about improving despite setbacks. \\
compassion & concern for others' suffering. & The author tells a story about a limping dog they saw on the sidewalk on their way home from work. \\
appreciation & recognition and enjoyment of the good qualities of person, place, or thing. & The author tells a story about their gratitude for a mailman who delivered mail in their neighborhood for years. \\
connection & closeness or shared understanding with a person, place, or thing. & The author tells a story about how they have been building new friendships and building a new home in a city they recently moved to.  \\
\midrule
\rowcolor{SkyBlue!15}
\multicolumn{3}{p{0.953\textwidth}}{
\textbf{Dimension: Character Appraisal}\newline
\textbf{Template}: \textit{While or after reading the story within this text, many readers from this subreddit would \{\{EITHER "positively", "negatively", or "neutrally"\}\} judge \{\{EITHER "narrator" OR IDENTIFIER/NAME OF OTHER CHARACTER FROM STORY\}\}.}\newline
\textbf{Summary}: Readers often judge whether character’s actions or state (e.g., beliefs, values, goals). Some common types of character appraisals include: positive, negative, or neutral judgment of the narrator (or another character).
} \\
\midrule
positive\_appraisal\_narr & a positive judgment of the narrator's actions or state. & In the story within the post, the author intervened to help someone who was being harassed at a bar. \\
negative\_appraisal\_narr & a negative judgment of the narrator's actions or state. & In the story within the post, the author made a snide comment to a colleague. \\
neutral\_appraisal\_narr & a neutral appraisal of the narrator's actions or state. & In the story within the post, the author describes a routine day in their life as a student. \\
positive\_appraisal\_other\_char & a positive judgment of a non-narrator character's actions or state. & In the story within the post, the author's partner gave them a gift. \\
negative\_appraisal\_other\_char & a negative judgment of a non-narrator character's actions or state. & In the story within the post, someone proctoring a test that the author was taking ignored that some people were cheating. \\
neutral\_appraisal\_other\_narr & a neutral appraisal of a non-narrator character's actions or state. & In the story within the post, someone ordered a type of coffee that the narrator was unfamiliar with. \\
\midrule
\rowcolor{SkyBlue!15}
\multicolumn{3}{p{0.953\textwidth}}{
\textbf{Dimension: Causal Explanation}\newline
\textbf{Template}: \textit{While or after reading the story within the post/comment, many readers from this subreddit would think that \{\{SHORT DESCRIPTION OF SITUATION/STATE/ACTION FROM STORY\}\} could be explained by \{\{SHORT EXPLANATION\}\}.}\newline
\textbf{Summary}: Readers often make inferences to explain aspects of a story. Some common types of explanatory inferences readers make include:\par - inferences about the narrator’s (or other characters’) values, beliefs, or goals (to explain character actions or state)\par - inferences about why an event or state not directly caused by a human happened (e.g., why a natural disaster happened)
} \\
\midrule
narr\_explained\_by\_narr & explaining some aspect of the narrator (e.g., their feelings or behavior) based on the perceived state (e.g., underlying values, beliefs, or goals) or actions of the narrator. & In the story within the post, the narrator quits their job and moves to a new city. \\
narr\_explained\_by\_other\_char\_or\_thing & explaining some aspect of the narrator (e.g., their feelings or behavior) based on (1) the perceived state (e.g., underlying values, beliefs, or goals) or actions of a non-narrator character or (2) the perceived state of affairs or event not directly attributed to any character. & In the story within the post, the narrator expresses their frustration with a conversation they had with a friend. \\
other\_char\_or\_thing\_explained\_by\_narr & explaining some aspect of (1) a non-narrator character (e.g., their feelings or behavior) or (2) the perceived state of affairs or event not directly attributed to any character based on the perceived state (e.g., underlying values, beliefs, or goals) or actions of the narrator. & In the story within the post, the narrator describes their depression and how it may be impacting their partner. \\
other\_char\_or\_thing\_explained\_by\_other\_char\_or\_thing & explaining some aspect of (1) a non-narrator character (e.g., their feelings or behavior) or (2) the perceived state of affairs or event not directly attributed to any character on (a) the perceived state (e.g., underlying values, beliefs, or goals) or actions of a non-narrator character or (b) some other perceived state of affairs or event not directly attributed to any character. & In the story within the post, the narrator describes how their manager has been more lax during the COVID-19 pandemic. \\
\midrule
\rowcolor{SkyBlue!15}
\multicolumn{3}{p{0.953\textwidth}}{
\textbf{Dimension: Prediction}\newline
\textbf{Template}: \textit{While or after reading the story within the post/comment, many readers from this subreddit would predict that \{\{EITHER "the narrator" or NAME/IDENTIFIER OF OTHER CHARACTER OR THING FROM STORY\}\} might \{\{SHORT DESCRIPTION OF ACTION OR STATE\}\}.}\newline
\textbf{Summary}: Readers often predict what will happen next in the world of a story. Some common types of predictions include:\par - predictions about the future actions or state of the narrator (or other characters) values, beliefs, or goals (to explain character actions or mental state)\par - predictions about a future event or state not directly caused by a human (e.g., an asteroid hitting)
} \\
\midrule
narr\_future\_state & a prediction about the future state of the narrator. & The story within the post ends with the narrator's partner breaking up with the narrator. \\
narr\_future\_action & a prediction about the future actions of the narrator. & The story within the post ends with the narrator's partner making an ultimatum for the narrator. \\
other\_char\_future\_state & a prediction about the future state of a non-narrator character. & The story within the post ends with the narrator breaking up with their partner. \\
other\_char\_future\_action & a prediction about the future actions of a non-narrator character. & The story within the post ends with the narrator making an ultimatum for their partner. \\
non\_char\_thing\_future\_state & a prediction about a future state not directly caused by a character  & The story within the post speculates about the state of technology in 100 years. \\
non\_char\_thing\_future\_event & a prediction about a future event not directly caused by a character  & The story within the post speculates on the impending demise of Black Friday sales due to increased online retail. \\
\midrule
\rowcolor{SkyBlue!15}
\multicolumn{3}{p{0.953\textwidth}}{
\textbf{Dimension: Stance}\newline
\textbf{Template}: \textit{After reading the story within this text, many readers from this subreddit would \{\{EITHER "support", "counter", or "be neutral to"\}\} the author's opinion\{\{SHORT DESCRIPTION OF AUTHOR'S OPINION/STANCE\}\}.}\newline
\textbf{Summary}: Readers often take a stance or overall opinion in response to a main idea, argument, or point advocated for by the author of a story, implicitly or explicitly. Readers can support (agree with), counter (disagree with), or be neutral to the author’s stance.
} \\
\midrule
support\_belief\_norm & a stance that mostly agrees with the stance most recently expressed (explicitly or implictly) in the preceding conversation.  & The story within the post is about how the narrator conducted extensive research to learn that the risk of shark attacks on the beaches of California is marginal. \\
counter\_belief\_norm & a stance that mostly disagrees with the stance most recently expressed (explicitly or implictly) in the preceding conversation.  & The story within the post is about how gambling is an effective way to get out of debt. \\
neutral\_belief\_norm & a stance that is neutral to the stance most recently expressed (explicitly or implictly) in the preceding conversation.  & The story within the post is about the narrator's strong opinion about a highly niche topic that most people don't know or care about. \\
\midrule
\rowcolor{SkyBlue!15}
\multicolumn{3}{p{0.953\textwidth}}{
\textbf{Dimension: Moral}\newline
\textbf{Template}: \textit{While or after reading the story within the post/comment, many readers from this subreddit would think that the moral of the story is \{\{MORAL/THEME\}\}.}\newline
\textbf{Summary}: Some common morals/values that stories highlight the importance of include: independence, novelty/change, pleasure, achievement, power, security, conformity, tradition, benevolence, and universalism.
} \\
\midrule
self-direction & independent thought and action, expressed in choosing, creating and exploring. & The story within the text is about how Copernicus discovered that the Earth revolves around the Sun despite the contemporary belief that the Earth was at the center of the solar system. \\
stimulation & excitement, novelty, and challenge in life. & The story within the text is about someone travelling internationally for the first time. \\
hedonism & pleasure, enjoyment, or sensuous gratification for oneself. & The story within the text is about how someone had an exhilarating experience seeing their favorite band at a concert.  \\
achievement & personal success through demonstrating competence according to social standards & The story within the text is about a person earning their nursing degree. \\
power & control or dominance over people and resources; social status and prestige & The story within the text is about how a new CEO changed the direction of a company and let go 200 employees. \\
security & safety, harmony, and stability of society, of relationships, and of self. & The story within the text is about how a narrator had a difficult conversation with their mother about their relationship dynamic which benefited it in the long run. \\
conformity & restraint of actions, inclinations, and impulses likely to upset or harm others and violate social expectations or norms in everyday interactions, usually with close others & The story within the text is about how the narrator navigated flirtatious conversations with a coworker when the narrator was in a committed relationship. \\
tradition & respect, commitment, and acceptance of the customs and ideas that one's culture or religion provides. & The story within the text is about how someone's religious beliefs and practices evolved over time. \\
benevolence & preserving and enhancing the welfare of those with whom one is in frequent personal contact (the 'in-group'). & The story within the text is about how the narrator's neighbors brought him food each night while he recovered from a surgery. \\
universalism & understanding, appreciation, tolerance, and protection for the welfare of all people and for nature. & The story within the text is about how people everywhere are struggling due to labor issues with the rise of automation. \\
\midrule
\rowcolor{SkyBlue!15}
\multicolumn{3}{p{0.953\textwidth}}{
\textbf{Dimension: Narrative Feeling}\newline
\textbf{Template}: \textit{While or after reading the story within the post/comment, the narrative content (i.e., the characters' situation and events) would spur many readers from this subreddit to feel \{\{FEELING/EMOTION\}\}.}\newline
\textbf{Summary}: Readers may sympathize or empathize with emotions of the narrator or other characters in a story, or they may feel emotions that differ from emotions of characters. The content of stories can also evoke memories in readers, causing readers to remember or relive emotions. Some of the feelings a reader could experience while or after reading a story include: fear, guilt, anger, sadness, disgust, envy, joy, pride, relief, hope, compassion, appreciation, and connection.
} \\
\midrule
fear & a response to perceived danger or threat. & The story within the post is about increased violence in their community. \\
guilt & remorse for violating personal or social standards. & The story within the post is about how the community (of which the reader is presumably a part) has failed to uphold their community standards. \\
anger & a strong reaction to perceived harm, injustice, or frustration. & The story within the post is an attack on a marginalized group. \\
sadness & a sense of loss, disappointment, or helplessness. & The story within the post is about the difficulty of making friends in a new city. \\
disgust & a reaction of revulsion to something perceived as offensive or repellent. & The story within the post is about a time the author accidentally sneezed into a customer's food while working at a restaurant and didn't tell them. \\
envy & a desire for something others have, coupled with resentment. & The story within the post is about the author winning the lottery. \\
joy & a state of happiness and contentment. & The story within the post depicts a heartwarming moment about two twins separated at birth reunited after 20 years. \\
pride & a sense of satisfaction from achievements or qualities. & The story within the post is about a community to which both the author and reader belong rebuilt itself after a terrible flood. \\
relief & a release from stress or tension after resolving a concern. & The story within the post is about how the author almost fell from a rock face during a free climb, but managed to save themselves. \\
hope & an optimistic expectation for a positive outcome. & The story within the post in a job forum is about how a company just posted 100 new job listings. \\
compassion & sympathy and concern for others' suffering. & The story within the post is about a fictional anthropomorphic mouse that got caught in a mouse trap. \\
appreciation & recognition and enjoyment of the good qualities of person, place, or thing. & The story within the post is about how the author has regularly picked up litter in a local creek for over a decade. \\
connection & closeness or shared understanding with a person, place, or thing. & The story within the post in a forum dedicated to a particular U.S. State recounts how the governor has stood up for the state in national politics. \\
\midrule
\rowcolor{SkyBlue!15}
\multicolumn{3}{p{0.953\textwidth}}{
\textbf{Dimension: Aesthetic Feeling}\newline
\textbf{Template}: \textit{While or after reading the story within the post/comment, the narrative form/techniques such as \{\{BRIEF DESCRIPTION OF TECHNIQUE OR FORMAL ELEMENT\}\} would spur many readers from this subreddit to feel \{\{FEELING\}\}.}\newline
\textbf{Summary}: The form, techniques, or style of a story can evoke aesthetic feelings in readers. Some aesthetic feelings a reader might experience include: suspense (anticipation for imminent event), curiosity (desire for information from the past to explain the present), surprise (experiencing a shocking/unexpected event), attentional engagement (feeling attention is grabbed and held by the story), feeling pulled into the world of the story, and visualization of images.
} \\
\midrule
suspense & a feeling of excitement or anxiety in anticipation of an imminent event. & The story within the text recounts the narrator's anxious week at work as they anticipated getting let go after hearing about company-wide layoffs.  \\
curiosity & desire for information from the past to explain the present & The story within the text introduces a character with a stab wound but is not told how they were wounded. \\
surprise & experiencing an unexpected or shocking event. & The story within the text describes the narrator accidentally cutting the tip of their finger off while making dinner. \\
attention\_engagement & finding the story to be compelling and capable of holding one's attention, as opposed to mundane or banal. & The story within the text describes a high-speed getaway from a bank robbery. \\
transportation & immersion, absorption, or feeling pulled into the world of the story. & The story within the text is suspenseful and describes how the narrator was running from someone chasing them late at night in the woods. \\
evocation & visualization, e.g. due to vivid language. & The story contains a whimsical description of the variously shaped clouds floating through the sky. \\
amusement & finding the story to be funny or to contain amusing elements. & The story within the text recounts the narrator mistaking their best friend for their romantic partner. \\
other & an aesthetic feeling not covered by the provided categories & \textasciitilde{} \\
\midrule
\end{longtable}
\endgroup